% ==========================================================================
% Moinul's LaTeX updates for the Group 24 report (29 Oct).
% Needs \usepackage{url} (already loaded if the report uses hyperref).
% Otherwise standard commands only: \texttt, tabular, itemize, enumerate, verbatim.
% Underscores inside \texttt are escaped as \_ .
% ==========================================================================

% --------------------------------------------------------------------------
% (1) Section 5.2: replace the first paragraph, ADD M6 after M5, and replace
%     the timing sentence at the end. M1 to M5 stay as they are.
% --------------------------------------------------------------------------

The method is implemented as a single Python 3.12 program organised in six
modules, run in order from one command (Appendix). It uses NumPy, OpenCV,
scikit-image, scikit-learn, PyTorch with torchvision, and matplotlib. Every
intermediate result is saved under the dataset folder
(\texttt{Dhan-Shomadhan/\_interim/}), so each module can be rerun on its own.

% ... M1 to M5 unchanged ...

\textbf{M6: Figures.} It draws Figures 1, 3 and 4 from the result files with
matplotlib, at full text width and at most 6~cm tall. The sample and failure
images are named in the program and read directly from the dataset folder, so
every figure can be rebuilt from the submission alone.

On a 16~GB M-series MacBook the modules took 46~s (M1), about 1~s (M2),
20~s (M3), 23~s (M4, Apple GPU) and 188~s (M5, all splits), and a few seconds
for M6: about 5 minutes in total.


% --------------------------------------------------------------------------
% (2) Appendix: replace ALL of Section 8 (parts a and b) with this.
%     Fits in 2 pages at 12 pt; check after compiling.
% --------------------------------------------------------------------------

\section{Appendix: User Manual}

\subsection*{a) Packages and Installation}

\begin{enumerate}
\item Install Python 3.12 from
\url{https://www.python.org/downloads/release/python-31210/} (3.12.10 is the
last 3.12 release with installers; our results were produced with 3.12.12). On
Windows, tick ``Add python.exe to PATH''.
\item Create an empty folder. Put \texttt{group24.py} (from
\texttt{group24.zip}) in it, and the unzipped dataset folder
\texttt{Dhan-Shomadhan} next to it, exactly as supplied (do not rename or move
its subfolders).
\end{enumerate}

\noindent Steps 3 to 7 are typed in a terminal:

\begin{center}
\small
\begin{tabular}{|>{\raggedright}p{2.6cm}|p{5.2cm}|p{5.6cm}|}
\hline
\textbf{Step} & \textbf{macOS Terminal} & \textbf{Windows Command Prompt} \\ \hline
3. Open the folder & \texttt{cd "/path/to/folder"} & \texttt{cd /d "C:\textbackslash folder"} \\ \hline
4. Create environment & \texttt{python3.12 -m venv venv} & \texttt{py -3.12 -m venv venv} \\ \hline
5. Activate it & \texttt{source venv/bin/activate} & \texttt{venv\textbackslash Scripts\textbackslash activate.bat} \\ \hline
6. Install packages & \multicolumn{2}{l|}{the command below (same on both)} \\ \hline
7. Check & \texttt{python -m pip check} & \texttt{python -m pip check} \\ \hline
\end{tabular}
\end{center}

\noindent Step 6, typed as ONE line (the exact versions used for our results):
{\small
\begin{verbatim}
python -m pip install numpy==2.5.4 matplotlib==3.11.2 opencv-python==4.13.0.92
  scikit-learn==1.9.1 scikit-image==0.26.0 torch==2.14.1 torchvision==0.29.1
\end{verbatim}}

\noindent Step 7 should print ``No broken requirements found''. torchvision is
part of the PyTorch project and provides the ResNet-18 model.

\paragraph{Pre-trained weights.} The ResNet-18 ImageNet weights
(\texttt{resnet18-f37072fd.pth}, about 45~MB) download automatically the first
time the code runs, so internet is needed for installation and the first run.
Without internet on that computer, download
\url{https://download.pytorch.org/models/resnet18-f37072fd.pth} elsewhere and
place it in \url{~/.cache/torch/hub/checkpoints/} (macOS/Linux) or
\url{%USERPROFILE%\.cache\torch\hub\checkpoints\} (Windows),
creating the folders if needed.

\subsection*{b) Running the Code to Reproduce the Results}

\begin{enumerate}
\setcounter{enumi}{7}
\item With the environment active and the terminal in the folder from step 2
(the folder containing \texttt{Dhan-Shomadhan}, as the brief's path
\texttt{.\textbackslash Dhan-Shomadhan\textbackslash} assumes), run:
\begin{verbatim}
python group24.py all
\end{verbatim}
\end{enumerate}

\noindent This runs six stages in order and prints progress; on a 16~GB
M-series MacBook it takes about 5 minutes. Each stage can also be run alone,
in this order, e.g.\ \texttt{python group24.py evaluate}. Later stages reuse
what earlier stages saved.

\begin{center}
\small
\begin{tabular}{|l|p{6.2cm}|p{5.2cm}|}
\hline
\textbf{Stage} & \textbf{What it does} & \textbf{Main output} \\ \hline
\texttt{inspect} & Checks labels, counts images, image statistics (Section 3) & \texttt{dataset-notes/} \\ \hline
\texttt{splits} & Makes the 5 seeded 80/20 train/test splits & \texttt{Dhan-Shomadhan/\_interim/} \\ \hline
\texttt{baseline} & Colour histogram and HOG features & \texttt{Dhan-Shomadhan/\_interim/} \\ \hline
\texttt{resnet} & ResNet-18 features (downloads weights on first use) & \texttt{Dhan-Shomadhan/\_interim/} \\ \hline
\texttt{evaluate} & Tunes and trains the SVMs, tests all splits and scenarios & \texttt{results/}, \texttt{results/logs/} \\ \hline
\texttt{plots} & Figures 1, 3 and 4 & \texttt{results/figures/} \\ \hline
\end{tabular}
\end{center}

\begin{center}
\small
\begin{tabular}{|p{5.6cm}|p{8.2cm}|}
\hline
\textbf{Report item} & \textbf{File} \\ \hline
Table 2 (mean $\pm$ std) and per-class recall & end of \texttt{results/logs/}\newline\texttt{05\_train\_evaluate\_log.txt} \\ \hline
Scores per split; per class & \texttt{results/summary.csv};\newline\texttt{results/per\_class.csv} \\ \hline
Confusion matrices; every test prediction & \texttt{results/confusion\_*.csv};\newline\texttt{results/predictions.csv} \\ \hline
Chosen SVM settings (Section 6.1) & \texttt{results/best\_params.csv} \\ \hline
Dataset statistics (Section 3) & \texttt{dataset-notes/} \\ \hline
\end{tabular}
\end{center}

\noindent\textbf{Reproducibility.} The splits use fixed seeds (0 to 4) and a
sorted file list, so they are identical on every computer. Deep features
computed on a different device (CPU instead of a GPU) may differ in the last
decimal places.

\noindent\textbf{Common problems.}
\begin{itemize}
\item \texttt{python3.12} or \texttt{py} not found: reopen the terminal after installing Python.
\item ``Dataset folder not found'': run from the folder that contains
\texttt{Dhan-Shomadhan}; check the unzip did not create a nested
\texttt{Dhan-Shomadhan/Dhan-Shomadhan/}.
\item ``No module named \ldots'': the environment is not active; repeat step 5.
\item Weight download fails: check the internet connection, or place the weights file manually (part a).
\end{itemize}


% --------------------------------------------------------------------------
% (3) Small accuracy fixes found in the full read-through (29 Oct)
% --------------------------------------------------------------------------

% (3a) Section 5.2: DELETE the red note "Moinul: please check this is still
%      accurate after your review".

% (3b) Algorithm 1, line 2: the code reports mismatches, it does not stop on
%      them (it stops only on an unknown folder name). Replace line 2 with:
%   Check each label against the CSV and the file-name prefix; report any mismatch

% (3c) Section 6.3, "Challenging Cases": the distance of field features from
%      white ones was not measured, so state it as an inference. Replace
%      "The field features lay far from all white training examples, so the RBF
%       SVM had no basis to separate them." with:
This suggests the field features lie far from all white training examples,
where an RBF SVM gives the same answer to every input.

% (3d) Section 6.1, "Hyperparameters" (the brief asks for ALL values): after
%      "RBF, C = 100 (4 splits) or 1000 (1 split) for the baseline" add
%      (check that page 8 still fits; if not, shorten the Table 2 caption):
, with gamma = $1/(4 \times 3844)$ (3 splits) or $1/(16 \times 3844)$ (2 splits)
